\section{Limitations and Conclusion}
\paragraph{Limitations.}
This is one training run with one seed, on 9$\times$9 Go, with a small network. We do not know how often
chained Elo misleads in other runs, games or network sizes; the analysis of noise is general, but the size
of the effect of shallow evaluation may not be. We compared only 32 and 64 simulations a move, and replayed only
generations 24--44. On a single pair, the log and the replay differ by less than the noise, so the attribution
rests on the aggregate. Wall-clock times were measured on a
machine shared with other work. We played no games against humans, and the GNU~Go comparison used a single
setting.

\paragraph{Conclusion.}
On a small AlphaZero run, summing Elo from short matches between consecutive generations suggested a
stall and decline. Anchored head-to-head games showed steady improvement. The disagreement is too large to
be match noise alone. Replaying the same consecutive matches with more search per move recovered the
correct direction and agreed with the anchored games. For small runs we recommend evaluating at least at
deployment search depth, and tracking scores against a few fixed anchors, which costs little. The training pipeline, the evaluation data and the identical-across-runtimes engine are
released with the paper \TODO{add the anonymized artifact link for review; the public repository URL for camera-ready}.

\paragraph{Use of AI tools.}
The implementation (training pipeline, rules engines, apps, evaluation scripts and
benchmarks) was written with an AI coding assistant (Claude Code) under the authors' direction and checked by
the tests and parity checks in Section~\ref{sec:results}. The same assistant, with a paper-writing skill,
drafted text, wrote the scripts that generate figures, tables and in-text numbers from the logged data, and
formatted and verified references. The authors reviewed all content and are responsible for it.
\TODO{author: confirm this statement matches what you did, and edit as needed.}
